\documentclass[12pt,a4paper]{article}

% ============================================================
% TEMPLATE - PROJETO DE ESTÁGIO SUPERVISIONADO
% Universidade Federal do Acre - Sistemas de Informação
% Baseado no modelo fornecido em DOCX.
% Compilador recomendado no Overleaf: pdfLaTeX
% ============================================================

\usepackage[utf8]{inputenc}
\usepackage[T1]{fontenc}
\usepackage[brazil]{babel}
\usepackage{graphicx}
\usepackage{geometry}
\usepackage{setspace}
\usepackage{indentfirst}
\usepackage{ragged2e}
\usepackage{titlesec}
\usepackage{fancyhdr}
\usepackage{tocloft}
\usepackage{microtype}
\usepackage{helvet}
\usepackage{url}
\usepackage[alf]{abntex2cite}

% Fonte sem serifa semelhante à Arial, compatível com pdfLaTeX.
\renewcommand{\familydefault}{\sfdefault}

% Margens do modelo original: superior/esquerda 3 cm; inferior/direita 2 cm.
\geometry{
  a4paper,
  top=3cm,
  bottom=2cm,
  left=3cm,
  right=2cm
}

% Espaçamento e parágrafos.
\onehalfspacing
\setlength{\parindent}{1.25cm}
\setlength{\parskip}{6pt}
\justifying

% ====================== DADOS EDITÁVEIS =====================
\newcommand{\Autor}{Kelvin Do Nascimento Moreira\\Marcos Antonio da Silva Manuares}
\newcommand{\TituloTrabalho}{Plataforma Operacional de Nivelamento, Triagem e Editais}
\newcommand{\Orientador}{Daricélio Moreira Soares}
\newcommand{\Curso}{CURSO DE BACHARELADO EM SISTEMAS DE INFORMAÇÃO}
\newcommand{\Instituicao}{UNIVERSIDADE FEDERAL DO ACRE}
\newcommand{\Cidade}{RIO BRANCO}
\newcommand{\Ano}{2026}
% ============================================================

% Formatação dos títulos.
\titleformat{\section}
  {\bfseries\normalsize}
  {\thesection}{0.5em}{\MakeUppercase}
\titlespacing*{\section}{0pt}{5cm}{3cm}

\titleformat{\subsection}
  {\bfseries\normalsize}
  {\thesubsection}{0.5em}{\MakeUppercase}
\titlespacing*{\subsection}{0pt}{1.9cm}{1.5cm}

% Comando para iniciar cada seção principal em uma nova página.
\newcommand{\secao}[1]{%
  \clearpage
  \section{#1}%
}

% Numeração de páginas no canto superior direito.
\pagestyle{fancy}
\fancyhf{}
\fancyhead[R]{\thepage}
\renewcommand{\headrulewidth}{0pt}
\setlength{\headheight}{15pt}

% Formatação do sumário.
\addto\captionsbrazil{\renewcommand{\contentsname}{SUMÁRIO}}
\renewcommand{\cfttoctitlefont}{\hfill\bfseries\normalsize}
\renewcommand{\cftaftertoctitle}{\hfill}
\renewcommand{\cftsecfont}{\bfseries}
\renewcommand{\cftsecpagefont}{\bfseries}
\renewcommand{\cftsubsecfont}{\bfseries}
\renewcommand{\cftsubsecpagefont}{\bfseries}
\renewcommand{\cftsecleader}{\cftdotfill{\cftdotsep}}
\renewcommand{\cftsubsecleader}{\cftdotfill{\cftdotsep}}
\setlength{\cftbeforetoctitleskip}{5cm}
\setlength{\cftaftertoctitleskip}{2.5cm}
\setlength{\cftbeforesecskip}{1pt}
\setlength{\cftbeforesubsecskip}{0pt}
\setlength{\cftsubsecindent}{0.95cm}
\setlength{\cftsecnumwidth}{0.8cm}
\setlength{\cftsubsecnumwidth}{1.2cm}

\begin{document}

% =========================== CAPA ============================
\begin{titlepage}
  \thispagestyle{empty}
  \centering

  \includegraphics[width=1.8cm]{ufac.png}\\[0.25cm]
  {\bfseries \Instituicao\par}
  {\bfseries \Curso\par}

  \vfill
  {\bfseries \TituloTrabalho\par}
  \vfill

  {\bfseries \Cidade\par}
  {\bfseries \Ano\par}
\end{titlepage}

% ===================== FOLHA DE ROSTO =======================
\begin{titlepage}
  \thispagestyle{empty}
  \centering

  {\bfseries \Autor\par}

  \vfill
  {\bfseries \TituloTrabalho\par}

  \vspace{3.2cm}
  \begin{flushright}
    \begin{minipage}{8.3cm}
      \setlength{\parindent}{0pt}
      Projeto de estágio supervisionado\\[0.35cm]
      Prof. Orientador: \Orientador
    \end{minipage}
  \end{flushright}

  \vfill
  {\bfseries \Cidade\par}
  {\bfseries \Ano\par}
\end{titlepage}

% ========================== SUMÁRIO ==========================
\clearpage
\pagestyle{empty}
\tableofcontents
\thispagestyle{empty}

% No modelo, a primeira seção aparece como página 3 no sumário.
\clearpage
\pagestyle{fancy}
\pagenumbering{arabic}
\setcounter{page}{3}

% ============================================================
% 1 APRESENTAÇÃO
% ============================================================
\section{APRESENTAÇÃO}

% Dar uma introdução ao tema. Falar sobre o que se trata o estágio,
% em qual empresa/setor será desenvolvido.
%
% Aproximadamente 1,5 página.

% Escreva o texto da apresentação abaixo deste comentário.

A transformação digital tem ampliado o uso de sistemas de software em diferentes áreas e aumentado a demanda por profissionais qualificados em Tecnologia da Informação. Esse cenário não ocorre de forma uniforme entre as regiões brasileiras. Na Região Norte, ainda existem desafios relacionados à oferta de formação especializada e à aproximação entre universidade e mercado de trabalho. Nesse contexto, iniciativas de formação complementar e parcerias entre instituições de ensino e empresas contribuem para ampliar a capacitação profissional e fortalecer o ecossistema regional de desenvolvimento de software \cite{vidal2026web,saar2023projeto}.

O Web Academy é uma iniciativa de extensão da Universidade Federal do Acre, realizada em parceria com a Motorola Mobility e a Flextronics, voltada à formação em desenvolvimento web \textit{full-stack}. O programa é desenvolvido no contexto da Lei de Informática e da Superintendência da Zona Franca de Manaus (SUFRAMA) e possui carga horária de 300 horas, distribuídas entre módulos básico, intermediário, avançado e \textit{hands-on}. A formação combina conteúdos técnicos, atividades práticas e desenvolvimento de projetos, com uso de metodologias ativas e práticas associadas a Scrum e Kanban \cite{saar2023projeto}.

Além da capacitação técnica, o Web Academy tem contribuído para o fortalecimento do ecossistema local de tecnologia, com reflexos na inserção de participantes no mercado de trabalho, no desenvolvimento de iniciativas empreendedoras e no fortalecimento da pesquisa e da pós-graduação em Computação \cite{vidal2026web}. Esse contexto cria também a necessidade de organizar de forma mais eficiente os processos administrativos relacionados à seleção e ao acompanhamento dos participantes.

O presente estágio está diretamente relacionado às atividades do Web Academy e tem como foco o desenvolvimento e a evolução do sistema P.O.N.T.E. (\textit{Plataforma Operacional de Nivelamento, Triagem e Editais}), criado para apoiar a execução dos processos seletivos do programa. A proposta surgiu da necessidade de reduzir atividades manuais e centralizar informações que antes dependiam de formulários, arquivos compactados, planilhas e diferentes procedimentos administrativos.

O P.O.N.T.E. reúne em uma única aplicação diferentes etapas do processo seletivo. Entre suas funcionalidades estão a criação e edição de editais, inscrição de candidatos, envio individual de documentos, triagem documental, análise curricular, recursos, registro de etapas secundárias, classificação, homologação, geração de documentos oficiais, comunicação com candidatos e exportação dos dados dos convocados para continuidade do processo no Portal do Aluno. O sistema também busca melhorar a rastreabilidade das decisões tomadas durante a seleção, principalmente por meio da associação dos indeferimentos aos critérios definidos no edital.

O desenvolvimento do sistema envolve conhecimentos de Engenharia de Software, desenvolvimento web, bancos de dados, modelagem de regras de negócio, segurança, testes e implantação. Dessa forma, o estágio permite aplicar conhecimentos adquiridos durante a graduação em uma situação real de desenvolvimento de software, acompanhando diferentes etapas do ciclo de vida da aplicação, desde a análise de requisitos até a implementação, testes, manutenção e implantação.

Assim, o estágio está inserido em um contexto de formação, desenvolvimento tecnológico e apoio às atividades do Web Academy. O trabalho com o P.O.N.T.E. busca contribuir para a organização e a automatização do processo seletivo, ao mesmo tempo em que proporciona experiência prática no desenvolvimento e na evolução de um sistema real.

% ============================================================
% 2 PROBLEMA DA PESQUISA
% ============================================================
\secao{PROBLEMA DA PESQUISA}

% Qual a grande questão que se busca responder?
% Qual o problema da empresa que se pretende resolver?
%
% Até uma página.

% Escreva o problema da pesquisa abaixo deste comentário.


% ============================================================
% 3 OBJETIVOS DA PESQUISA
% ============================================================
\secao{OBJETIVOS DA PESQUISA}

\subsection{OBJETIVO GERAL}

% Qual o objetivo maior do trabalho?

% Escreva o objetivo geral abaixo deste comentário.


\subsection{OBJETIVOS ESPECÍFICOS}

% Quais os demais objetivos?

% Escreva os objetivos específicos abaixo deste comentário.


% ============================================================
% 4 JUSTIFICATIVA DA PESQUISA
% ============================================================
\secao{JUSTIFICATIVA DA PESQUISA}

% Por que foi realizado o estágio sobre este assunto, nesta empresa?
% Justificar aqui a importância do tema.
% Qual a importância em se pesquisar este assunto?
%
% Aproximadamente 1 página.

% Escreva a justificativa abaixo deste comentário.


% ============================================================
% 5 FUNDAMENTAÇÃO TEÓRICA
% ============================================================
\secao{FUNDAMENTAÇÃO TEÓRICA}

% Parte teórica do estágio: qual o assunto que dá o embasamento para o estágio?
% Quais os principais temas que precisam ser abordados?
%
% Geralmente de 3 a 5 páginas baseadas em autores, de preferência livros,
% usando as normas metodológicas para citá-los.

\subsection{SUBTÍTULO}

% Escreva aqui o conteúdo do primeiro subtópico.


\subsection{SUBTÍTULO}

% Escreva aqui o conteúdo do segundo subtópico.


\clearpage
\subsection{SUBTÍTULO}

% Escreva aqui o conteúdo do terceiro subtópico.


\subsection{SUBTÍTULO}

% Escreva aqui o conteúdo do quarto subtópico.


% ============================================================
% 6 PROCEDIMENTOS METODOLÓGICOS
% ============================================================
\secao{PROCEDIMENTOS METODOLÓGICOS}

% Como será realizado o estágio? Pesquisa bibliográfica, estudo de caso,
% levantamento (ver manual sobre metodologia no grupo).
%
% Detalhar se há instrumentos de coleta de dados, como questionário,
% formulário, entrevista ou observação; quem é a população e a amostra;
% e como os dados serão tabulados.
%
% Se for desenvolvimento de sistemas ou de aplicações, falar quais
% ferramentas serão utilizadas.

% Escreva os procedimentos metodológicos abaixo deste comentário.


% ============================================================
% 7 ESBOÇO DOS CAPÍTULOS E SEÇÕES
% ============================================================
\secao{ESBOÇO DOS CAPÍTULOS E SEÇÕES}

% O relatório de estágio seguirá a seguinte estrutura:
%
% Detalhar a estrutura que será adotada. Sugere-se consultar o modelo
% de relatório de estágio e detalhar principalmente a parte do referencial
% teórico, que mudará para cada caso.

% Escreva o esboço dos capítulos e seções abaixo deste comentário.


% ============================================================
% 8 CRONOGRAMA
% ============================================================
\secao{CRONOGRAMA}

% Colocar um gráfico de Gantt com a programação do estágio, considerando
% prazos para pesquisas bibliográficas, elaboração do projeto, apresentação
% do projeto, desenvolvimento do estágio na empresa, elaboração/aplicação
% de questionários, tabulação de dados, entrevistas, elaboração de
% relatórios etc.
%
% Você pode inserir aqui uma tabela em LaTeX ou uma imagem do gráfico de Gantt.

% Insira o cronograma abaixo deste comentário.


% ============================================================
% REFERÊNCIAS BIBLIOGRÁFICAS
% ============================================================
\clearpage
\renewcommand{\refname}{REFERÊNCIAS BIBLIOGRÁFICAS}
\phantomsection
\addcontentsline{toc}{section}{REFERÊNCIAS BIBLIOGRÁFICAS}

% As referências devem ser cadastradas no arquivo referencias.bib.
%
% Para citar no texto, use, por exemplo:
% \citeonline{chave}  -> Autor (ano)
% \cite{chave}        -> (AUTOR, ano)

\bibliography{referencias}

\end{document}
